\documentclass[11pt]{article}
\usepackage{amsmath,amssymb,amsthm}
\newtheorem{theorem}{Theorem}
\newtheorem{lemma}[theorem]{Lemma}
\title{A Linear Bound for a Common Moving Pole}
\author{}
\date{November 23, 2026}
\begin{document}
\maketitle

\paragraph{Status and scope.}
This note gives a complete written argument for a restricted readout
family. It is not Lean checked. It does not bound the original model,
in which different heads have independent denominator parameters.
In particular it proves neither accepted global target. The purpose
is to decide whether a proposed one-parameter Cauchy-plane sweep can
supply the missing superconstant local saving. It cannot.

Fix $H\geq1$, an interval $I\subset\mathbb R$, and, for each query
index $i$, real numbers $p_{ih},A_{ih},c_i$. Assume
$p_{ih}-t>0$ for every $t\in I$ and every $i,h$. Consider the score
family, with parameters $t\in I$ and $a\in\mathbb R^H$,
\[
 F_i(t,a)=c_i+\sum_{h=1}^H\frac{A_{ih}+a_h}{p_{ih}-t}.
 \tag{1}
\]
The query poles and numerator offsets may vary independently with
both $i$ and $h$. Only the moving pole parameter $t$ is common.

\begin{theorem}
The strict-sign class in (1) cannot shatter $64(H+1)$ query indices.
Consequently its VC dimension is less than $64(H+1)$, uniformly over
all choices of query data and of the positive-pole interval $I$.
\end{theorem}

We use two elementary finite-dimensional facts, with their needed
justifications recorded here.

\begin{lemma}[Minor signs determine strict feasibility]
Let $v_1,\ldots,v_m\in\mathbb R^d$. The set of sign vectors
$s\in\{-1,1\}^m$ for which some $w\in\mathbb R^d$ satisfies
$s_i\langle v_i,w\rangle>0$ for every $i$ is determined by the signs
of all square minors of the matrix with rows $v_i$.
\end{lemma}
\begin{proof}
For a fixed $s$, the strict theorem of alternatives says that such
$w$ exists exactly when there is no nonzero vector $q\geq0$ with
$\sum_iq_i s_iv_i=0$. This alternative follows, for example, by
strictly separating zero from the compact convex hull of the
$s_iv_i$: a nonnegative dependence can be normalized to a convex
combination, and strict separation gives the required $w$.

Any nonnegative dependence has a support-minimal one. Its support
is a minimally dependent set: if the kernel on that support had
dimension at least two, one could perturb the strictly positive
coefficients within the kernel until one coefficient became zero,
contradicting minimality. A minimally dependent support of size
$k+1$ has rank $k\leq d$. Choose $k$ independent coordinate columns
on its rows. Its one-dimensional dependence has coefficients given,
up to a common nonzero scalar, by the alternating $k$-row minors.
Their signs determine whether multiplication by $s_i$ makes every
nonzero coefficient have the same sign. A zero row, the case $k=0$,
is also detected by the one-entry minors. Thus all possible positive
minimal dependencies, and hence strict feasibility, are determined
by the signs of the square minors.
\end{proof}

\begin{lemma}[Hyperplane sign count]
For $m$ vectors in $\mathbb R^d$, at most
$2\sum_{j=0}^{d-1}\binom{m-1}{j}$ strict sign vectors are possible.
\end{lemma}
\begin{proof}
The complement of the $m$ homogeneous hyperplanes has one convex
open chamber for each realized strict sign vector. Adding the last
hyperplane creates at most as many new chambers as the preceding
hyperplanes induce on that hyperplane, a space of dimension $d-1$.
Thus the maximal chamber counts obey
$R(m,d)\leq R(m-1,d)+R(m-1,d-1)$, with $R(1,d)\leq2$ and
$R(m,1)\leq2$. Induction and Pascal's identity give the bound.
Zero or repeated hyperplanes can only decrease the count.
\end{proof}

\begin{proof}[Proof of the theorem]
Put $d=H+1$ and consider any $r$ query indices. Define the $r+1$
row vectors
\[
 v_i(t)=\left(\frac1{p_{i1}-t},\ldots,
          \frac1{p_{iH}-t},\ c_i+\sum_h\frac{A_{ih}}{p_{ih}-t}\right),
 \qquad v_{r+1}(t)=(0,\ldots,0,1).
\]
A strict sign vector $s$ of (1) is feasible exactly when the
augmented homogeneous sign vector $(s,+1)$ is feasible for these
rows. Indeed, the final inequality makes the last coordinate of
the separating vector positive; division by that coordinate
normalizes it to $(a,1)$.

There are at most
\[
 K=\sum_{k=0}^{d}\binom{r+1}{k}\binom{d}{k}
   =\binom{r+d+1}{d}
\]
square minors, counting also the empty minor. In a minor using $k$
ordinary query rows, multiply by the product of their $kH$
positive affine denominators. Each resulting matrix row has
polynomial entries of degree at most $H$. Multilinearity of the
determinant therefore gives a numerator of degree at most
$kH\leq d^2$. The normalization row introduces no denominator.
Identically zero numerators have constant zero sign and are omitted.
The remaining minor numerators have at most $d^2K$ distinct real
zeros altogether. These zeros split $I$ into at most $2d^2K+1$
pieces, counting singleton zeros as pieces. All minor signs are
constant on each open piece. By the first lemma the feasible sign
vectors are constant there as well. Singleton pieces are handled
by their actual matrices, including any rank drops.

The second lemma now bounds the total number of strict sign
vectors in (1) by
\[
 (2d^2K+1)\,2\sum_{j=0}^{d-1}\binom rj.
 \tag{2}
\]
There is no genericity assumption in this count.

Set $r=64d$. The elementary binomial bounds
$\binom Nk\leq(eN/k)^k$ and
$\sum_{j=0}^k\binom Nj\leq(eN/k)^k$ for $1\leq k\leq N/2$
give, using $e<3$ and $d\geq2$,
\[
 K\leq(66e)^d<2^{8d},\qquad
 \sum_{j=0}^{d-1}\binom rj\leq(128e)^{d-1}<2^{9d}.
\]
For clarity, the sum inequality follows by taking $x=k/N\leq1$:
$x^k\sum_{j=0}^k\binom Nj\leq(1+x)^N\leq e^k$.
Also $d^2\leq2^{2d}$ and $2d^2K+1\leq4d^2K$. Therefore (2) is
less than $2^{19d+3}$, which is less than $2^{64d}$. These $r$
queries are not shattered. Any larger shattered set would contain
a shattered subset of this size, proving the assertion.
\end{proof}

\paragraph{What remains available.}
The argument allows arbitrary row poles, numerator offsets, row
biases, and residue sizes. It also allows all cancellations and
all singular parameter values. It does not apply to independently
moving poles $u_h$, nor to the arbitrary-permutation planes with
independent token assignments. The degree factor here contributes
only a polynomial factor to the number of parameter intervals;
it cannot be raised to $H$ merely because there are $H$ residue
parameters. Several moving pole parameters require a separate
multivariate count and do not fall under the theorem as stated.
\end{document}
